\section{Related work}
\label{sec:related_work}

\textbf{Knowledge Graph Question Answering (KGQA).}
To address the limitations of unstructured text retrieval and LLM hallucination, KGs have been widely adopted in question answering tasks for their ability to provide structured world knowledge \citep{baek2023knowledgeaugmentedlanguagemodelverification, xu2025knowledgeinfusedpromptingassessingadvancing}.
Traditional KGQA methods fall into two categories: \textbf{Information retrieval (IR)} methods \citep{Zhang_2022, dong2023bridgingkbtextgapleveraging, li2024chainofknowledgegroundinglargelanguage} retrieve subgraphs from KGs and extract answers from the gathered evidence, while \textbf{Semantic parsing (SP)} methods \citep{berant2013semantic, yih2015semantic, li2023fewshotincontextlearningknowledge, Luo_2024} translate questions into logical forms (e.g., SPARQL) for verifiable execution.
However, both are brittle under noisy or incomplete KG schemas and struggle to handle long-tail entities, motivating the shift toward agentic paradigms.

\textbf{Agentic KGQA.}
Agentic KGQA treats LLMs as autonomous agents that interact with 
KGs through iterative exploration and retrieval.
\textbf{Prompting-based} methods such as ToG \citep{sun2023think} and GoG \citep{xu2024generate} perform iterative graph traversal but rely on predefined workflows or costly closed-source LLMs.
\textbf{Fine-tuning} methods learn from curated training trajectories: CoT-based approaches \citep{wu2025cotkrchainofthoughtenhancedknowledge, Zhao2024KGCoTCP} enhance step-by-step reasoning transparency, RoG \citep{luo2023reasoning} grounds plans in KGs, and KG-Agent \citep{jiang2025kg} employs multi-agent reasoning but 
remains limited to synthesized program data.
\textbf{Reinforcement learning} methods such as KG-R1 \citep{song2025efficient} optimize the retrieval policy end-to-end but still rely on subgraph pre-extraction.

